% Slajd 10 – zespół, prośba do jury i partnerów, zamknięcie. Wiersze zespołu ustawia \TeamRows w main.tex.
\begin{frame}{Zespół i~czego potrzebujemy}
  \begin{columns}[T,onlytextwidth]
    \begin{column}{0.46\textwidth}
      \begin{block}{Zespół}
        \footnotesize
        \renewcommand{\arraystretch}{1.2}
        \begin{tabular}{@{}>{\bfseries}l l@{}}
          \TeamRows
        \end{tabular}
        \vspace{1mm}

        Prototyp zbudowany podczas HackYeah; kod w~otwartym repozytorium, gotowy do uruchomienia jednym poleceniem.
      \end{block}
    \end{column}
    \begin{column}{0.52\textwidth}
      \begin{block}{Czego szukamy po hackathonie}
        \footnotesize
        \begin{itemize}
          \item \textbf{3~lokale i~1~instytucja kultury} w~Krakowie do 30-dniowego pilotażu panelu właściciela.
          \item \textbf{Partner społeczny} (organizacja osób z~niepełnosprawnościami) do weryfikacji danych
                i~testów dostępności z~użytkownikami.
          \item \textbf{Kontakt do otwartych danych drugiego miasta} i~do organizacji turystycznych jako kanału.
          \item \textbf{Finansowanie przedwdrożeniowe} na 6~miesięcy (akcelerator, grant na dostępność),
                zanim przychody z~planu Pro pokryją zespół.
        \end{itemize}
      \end{block}
    \end{column}
  \end{columns}
  \vspace{2.5mm}
  \centering
  {\bfseries\color{accNavy}Nie kolejna statyczna mapa -- żywy, weryfikowalny obraz dostępności miasta, który sam na siebie zarabia.\par}
  \vspace{0.8mm}
  \muted{Demo: \DemoUrl\ \textperiodcentered\ Kod: \RepoUrl\ \textperiodcentered\ Film: \VideoUrl}

  \note[item]{Zakończyć zdaniem z~tego slajdu i~zaprosić do stolika z~telefonem -- niech jury samo kliknie „Zgłoś problem”.}
  \note[item]{Przy prośbie o~finansowanie podać rząd wielkości, jeśli jury zapyta: 6~miesięcy pracy 2--3 osób plus infrastruktura; przychody z~wariantu bazowego pojawiają się od drugiego roku.}
  \note[item]{Pytanie o~zespół po hackathonie: operator (spółka / fundacja), moderacja w~panelu, finansowanie z~planu Pro i~wdrożeń white-label.}
\end{frame}
